%%%
% Run this file multiple times to ensure TOC, Bibliography, References, etc are populated correctly
%
% For font and typeface related issues, edit the book-template.sty file
%

%%%
% Preamble
%

\documentclass{sumintegrals-book}

%%%
% Be aware there are issues with bold math, so \mathbf{...} will work but not in section heads or chapter titles.
%
% To fix some issues, you may want to consider using the following package:
%\usepackage{mathptmx}
%
% For postscript text:
\usepackage{sumintegrals-book}
\usepackage{imakeidx}

%%%
% Some other packages to consider are:
%
%\usepackage{lipsum} % This will make "junk" text
\usepackage{chapterbib} % This will make chapter bibliographies with BibTeX
%\usepackage{multind} % This is useful for multiple indices
%\usepackage{answers} % This is used for answers to problems

%%%
% Set the depth of section head numbering for Table of Contents:
% 0=Chapters, 1=Sections, 2=Subsections, 3=Subsubsections
%
\setcounter{tocdepth}{1}
\setcounter{chapter}{0} % This sets the "first" chapter number.

%%%
% To make the document visibly a "draft" document, use the following command which adds double spacing between lines, current date/time at the footer.
%\draft
%
% To prevent tables from double spacing with this activated, use
%\renewcommand{\arraystretch}{0.6}

%%%
% Greek Alphabet
% The letters are not all defined, so here are the commands added to ensure all Greek letters appear:
\newcommand{\Alpha}{\mathrm{A}}
\newcommand{\Beta}{\mathrm{B}}
\newcommand{\Epsilon}{\mathrm{E}}
\newcommand{\Zeta}{\mathrm{Z}}
\newcommand{\Eta}{\mathrm{H}}
\newcommand{\Iota}{\mathrm{I}}
\newcommand{\Kappaa}{\mathrm{K}}
\newcommand{\Mu}{\mathrm{M}}
\newcommand{\Nu}{\mathrm{N}}
\newcommand{\Omicron}{\mathrm{O}}
\newcommand{\Rho}{\mathrm{P}}
\newcommand{\Tau}{\mathrm{T}}
\newcommand{\Chi}{\mathrm{X}}

%%%
% Generic Quantifier Q
\newcommand{\quantif}{\hbox{\tikz{
\draw[line width=0.45pt] (0,0) circle (3pt and 4pt);
\draw[line width=0.45pt, line cap=round] (1pt,-2pt) -- ++(2pt,-2pt);
}}}

%%%
% Put the index in a Table of Contents
\makeindex[intoc]

% END OF PREAMBLE
%%%
% BEGIN OF DOCUMENT
%

\begin{document}

%%%
% Set the main title, subtitle, author(s), affiliation(s), etc.
\booktitle{Discrete Mathematics}
\subtitle{with Logic and Set Theory}

\authors{Sum Integrals\\
\affil{$\Sigma\int s$}
}

\offprintinfo{Discrete Mathematics}{$\Sigma\int s$}
% \\ Can be used if the title, edition, etc are too wide

%%%
% FRONT COVER (sometimes omitted for PDF)
%\frontcoverpage

%%%
% FRONT PAGE
\halftitlepage

%%%
% INSIDE BOOK FRONT PAGE
\titlepage

%%%
% DEDICATION PAGE
\dedication{For all those who helped me along the way.}

%%%
% CONTRIBUTORS
%\begin{contributors}
%\name{First Last}, Affiliation, Location.
%\name{First Last}, Affiliation, Location.
%\end{contributors}

%%%
% CONTENTS IN BRIEF AND TABLE OF CONTENTS
%\contentsinbrief
\tableofcontents

%%%
% LIST OF FIGURES AND LIST OF TABLES
%\listoffigures
%\listoftables

%%%
% FORWARD
%\begin{forward}
%Forward text.
%\end{forward}

%%%
% SECTION/CHAPTER DIAGRAM
$$
\begin{tikzcd}
\textbf{\thetitle}&\\
\boxed{\bf CH1}\arrow[d]&\\
\boxed{\bf CH2}\arrow[d]\arrow[r]&\boxed{\bf CH3}\\
\boxed{\bf CH4}\arrow[d]&\\
\boxed{\bf CH5}\arrow[d]&\\
\boxed{\bf CH6}\arrow[d]\arrow[r]&\boxed{\bf CH7}\arrow[d]\\
\boxed{\bf CH9}&\boxed{\bf CH8}
\end{tikzcd}
$$

%%%
% PREFACE
\begin{preface}
Preface text.
\end{preface}

%%%
% ACKNOWLEDGEMENTS
%\begin{acknowledgements}
%Acknowledgements text.
%\end{acknowledgements}

%%%
% ACRONYM LIST
\begin{acronyms}
\acro{WLOG}{Without Loss of Generality}
\acro{ISTS}{It Suffices to Show}
\acro{TFAE}{The Following Are Equivalent}
\acro{AFC}{Assume for Contradiction}
\acro{PLAE}{Proof Left as Exercise}
\end{acronyms}

%%%
% GLOSSARY LIST
%\begin{glossary}
%\term{theTerm}Description of the term.
%
%\term{theTerm}Description of the term.
%
%\end{glossary}

%%%
% SYMBOL LIST
\begin{symbols}
\fbunch{Greek Alphabet} % FIRST BUNCH
\term{\Alpha~\alpha}{Alpha}
\term{\Beta~\beta}{Beta}
\term{\Gamma~\gamma}{Gamma}
\term{\Delta~\delta}{Delta}
\term{\Epsilon~\e}{Epsilon}
\term{\Zeta~\zeta}{Zeta}
\term{\Eta~\eta}{Eta}
\term{\Theta~\theta~\vartheta}{Theta}
\term{\Iota~\iota}{Iota}
\term{\Kappa~\kappa~\varkappa}{Kappa}
\term{\Lambda~\lambda}{Lambda}
\term{\Mu~\mu}{Mu}
\term{\Nu~\nu}{Nu}
\term{\Xi~\xi}{Xi}
\term{\Omicron~o}{Omicron}
\term{\Pi~\pi~\varpi}{Pi}
\term{\Rho~\rho}{Rho}
\term{\Sigma~\sigma~\varsigma}{Sigma}
\term{\Tau~\tau}
\term{\Upsilon~\upsilon}{Upsilon}
\term{\Phi~\phi~\varphi}{Phi}
\term{\Chi~\chi}{Chi}
\term{\Psi~\psi}{Psi}
\term{\Omega~\omega}{Omega}

\bunch{Propositional Logic} % BUNCH
\term{\wedge\bigwedge_{i}}{Conjunction}
\term{\vee\bigvee_{i}}{Disjunction}
\term{\neg}{Negation}
\term{\implies}{(Material) Conditional}
\term{\iff}{(Material) Biconditional}
\term{oplus}{Exclusive Disjunction}
\term{textbf{T},\texttt{1}}{Truth-Value True}
\term{\textbf{F},\texttt{0}}{Truth-Value False}
\term{\top}{Tautology}
\term{\bot}{Contradiction}
\term{\therefore}{Argument ``Therefore''}
\term{\vDash}{Logical Entailment}
\term{\vdash}{Logical Derivability}
\term{\equiv}{Logical Equivalence}

\bunch{Set Theory} % BUNCH
\term{\cup\bigcup_{i}}{Union}
\term{\cap\bigcap_{i}}{Intersection}
\term{\overline{A},A^{\complement}}{Complement}
\term{A-B,A\setminus B}{Difference}
\term{\powerset(A)}{Power Set}
\term{A\times B}{Cartesian Product}
\term{A^{n}}{Product Space}
\term{\in}{Element of}
\term{\subseteq}{Subset of}
\term{subset}{Proper Subset of}
\term{=}{Equality}
\term{\#A}{Set Cardinality}
\term{\abv{A}}{Set Ordinality}
\term{\varnothing}{Null Set}
\term{\nats}{Natural Numbers $\{1,2,\ldots\}$}
\term{\nats_{0}}{Whole Numbers $\{0,1,2,\ldots\}$}
\term{\ints}{Integer Numbers}
\term{\rats}{Rational Numbers}
\term{\algs}{Algebraic Numbers}
\term{\reals}{Real Numbers}
\term{\complex}{Complex Numbers}
\term{\aleph_{0}}{Countable Infinity}
\term{\mathfrak{c}}{Continuum Cardinality}

\bunch{Predicate Logic} % BUNCH
\term{\forall}{Universal Quantifier}
\term{\exists}{Existential Quantifier}
\term{\quantif}{Generic Quantifier}
\term{\exists!}{Uniqueness Quantifier}

\bunch{Functions} % BUNCH
\term{f:A\to B}{Function from $A$ to $B$}
\term{a\longmapsto b}{Maps $a$ to $b$}
\term{\mathrm{im}(f)}{Image of $f$}
\term{f(A)}{Range of $f$ over $A$}
\term{f^{-1}}{Preimage and Inverse of $f$}
\term{\mathrm{id}_{A}}{Identity Function on $A$}
\term{f\big|_{A}}{$f$ Restricted to $A$}
\term{A^{B}}{Set of Functions from $A$ to $B$}
\term{\hookrightarrow}{Injective Map}
\term{??}{Surjective Map}

\bunch{Arithmetic} % BUNCH
\term{+,\Sigma}{Addition,Sum}
\term{-}{Subtraction,Difference}
\term{\cdot,\Pi}{Multiplication,Product}
\term{/}{Division,Quotient}

\bunch{Number Theory} % BUNCH
\term{a\,\big|\,b}{$a$ Divides $b$}
\term{a~\mathbf{div}~m}{Division Modulo $m$}
\term{a~\mathbf{mod}~m}{Remainder Modulo $m$}
\term{\equiv}{Congruence}
\term{(\mathrm{mod}\,m)}{Modulo $m$}
\term{\mathrm{gcd}(a,b)}{Greatest Common Divisor}
\term{\mathrm{lcm}(a,b)}{Least Common Multiple}


\end{symbols}
